
sur @@M0@@, d'où le
résultat, l'homomorphisme étant compatible aux filtrations.

\section*{6. Anneau de Chern}

Dans ce paragraphe, l'anneau de base sera un anneau binomial (2.7) @@M1@@, fixé une fois pour toutes. Nous utiliserons la catégorie @@M2@@ (1.1), et par suite reviendrons aux notations du paragraphe 1, réservant la notation @@M3@@ au foncteur défini sur @@M4@@ (1.1), et notant @@M5@@ sa restriction à @@M6@@ (1.8).

6.1. Nous avons défini en 3.5 et 3.6 des morphismes de foncteurs :
@@M7@@
le dernier étant la multiplication pour la loi @@M8@@ par @@M9@@, avec @@M10@@.

Par 1.9, on en déduit des morphismes :
@@M11@@
et il résulte des propriétés d'isobarité des polynômes @@M12@@, @@M13@@, @@M14@@ (3.5, 3.6, 2.9) que l'image du premier et du troisième foncteurs est dans @@M15@@, et celle du deuxième, dans le sous-groupe des séries formelles en @@M16@@ dont les coefficients en degré non nul sont dans @@M17@@.

En composant le premier et le troisième avec le morphisme canonique : @@M18@@ (1.11), on obtient d'après 1.12 deux morphismes de
foncteurs sur @@M19@@ :
